\documentclass[10pt]{beamer}

\usetheme{metropolis}
\usepackage{appendixnumberbeamer}
\usepackage{booktabs}
\usepackage{listings}
\usepackage{xcolor}
\usepackage{colortbl}
\usepackage{array}

% Python listing style
\definecolor{codebg}{HTML}{F5F5F5}
\definecolor{codegreen}{HTML}{28A745}
\definecolor{codepurple}{HTML}{6F42C1}
\definecolor{codeorange}{HTML}{D73A49}
\definecolor{codegray}{HTML}{6A737D}
\lstset{
  language=Python,
  basicstyle=\ttfamily\tiny,
  keywordstyle=\color{codeorange}\bfseries,
  commentstyle=\color{codegray}\itshape,
  stringstyle=\color{codegreen},
  backgroundcolor=\color{codebg},
  frame=single,
  framerule=0.4pt,
  breaklines=true,
  showstringspaces=false,
  columns=fullflexible,
  keepspaces=true,
  tabsize=4,
  xleftmargin=4pt,
  xrightmargin=4pt,
  aboveskip=4pt,
  belowskip=4pt,
  morekeywords={jnp,lax,self,replace},
}

% Green highlight for table rows
\definecolor{rowgreen}{HTML}{D4EDDA}

\title{Porting DFlash Speculative Decoding\\from GPU/PyTorch to TPU/JAX}
\subtitle{UCSD DSC 180 Capstone}
\date{Winter 2025}
\author{Aaron Feng \and Zhongyan Luo \and Son Nguyen \and Andy Huang}
\institute{Mentor: Hao Zhang \quad TA: Yiming Zhao}

\begin{document}

% ============================================================
% Slide 1 -- Title
% ============================================================
\maketitle

% ============================================================
% Slide 2 -- Problem: The Autoregressive Bottleneck
% ============================================================
\section{Problem}

\begin{frame}{The Autoregressive Bottleneck}
\begin{columns}[T]
\begin{column}{0.55\textwidth}
\begin{itemize}
  \item Standard LLM decoding: \textbf{1 forward pass per token}
  \item Latency scales linearly with output length
  \item \textbf{Speculative decoding:} cheap draft model proposes tokens, target model verifies in parallel
  \item Key metric: \textbf{$\tau$} = expected accepted tokens per step
  \item Existing drafters (Eagle, Medusa) cap at $\tau \approx 2$--$3$
\end{itemize}
\end{column}
\begin{column}{0.42\textwidth}
\vspace{1em}
\centering
\footnotesize
\begin{tabular}{lc}
\toprule
\textbf{Method} & \textbf{Fwd passes} \\
\midrule
Autoregressive & $n$ \\
Spec.\ Decoding & $\sim n / \tau$ \\
\bottomrule
\end{tabular}

\vspace{1em}
\footnotesize
Higher $\tau$ $\Rightarrow$ fewer passes $\Rightarrow$ lower latency.

\vspace{0.5em}
With $\tau = 7$: generate 7 tokens per step instead of 1.
\end{column}
\end{columns}
\end{frame}

% ============================================================
% Slide 3 -- DFlash: Block Diffusion Drafting
% ============================================================
\begin{frame}{DFlash: Block Diffusion Drafting}
\begin{itemize}
  \item Drafts \textbf{16 tokens in one forward pass} (block diffusion, non-causal attention)
  \item Conditioned on target hidden states from layers \texttt{[1, 9, 17, 25, 33]}
  \item Shares target model's embedding layer and LM head (no extra parameters)
  \item GPU baseline: $\tau = 7.07$, \textbf{5.53$\times$ speedup}
\end{itemize}

\begin{alertblock}{Why DFlash is different}
1 draft pass (16 tokens) + 1 verify pass \quad vs.\quad Eagle3: 3 sequential drafts + 1 verify
\end{alertblock}
\end{frame}

% ============================================================
% Slide 4 -- Our Contribution
% ============================================================
\section{Approach}

\begin{frame}{Our Contribution}
\begin{itemize}
  \item \textbf{Goal:} port DFlash from GPU/PyTorch $\rightarrow$ TPU/JAX in \texttt{tpu-inference}
  \item Target: \textbf{Qwen3-4B} \quad Draft: \textbf{DFlash-b16}
  \item Full \textbf{vLLM integration} --- zero vLLM source changes required
  \item Standalone benchmark harness for direct GPU comparison
  \item Hardware: \textbf{TPU v4-8} (4 chips, 8 cores)
\end{itemize}
\end{frame}

% ============================================================
% Slide 5 -- Architecture: Dual KV Cache
% ============================================================
\begin{frame}{Architecture: Dual KV Cache}
\begin{columns}[T]
\begin{column}{0.48\textwidth}
\textbf{Target Model}
\begin{itemize}
  \item Paged KV cache
  \item \texttt{ragged\_paged\_attention}
  \item Extracts \texttt{aux\_hidden\_states} from layers [1,9,17,25,33]
\end{itemize}
\end{column}
\begin{column}{0.48\textwidth}
\textbf{Draft Model (DFlash)}
\begin{itemize}
  \item Static KV cache
  \item \texttt{dynamic\_update\_slice} writes
  \item \texttt{flash\_attention} (causal=False)
\end{itemize}
\end{column}
\end{columns}

\vspace{1em}
\centering
\small
Target aux features $\xrightarrow{\text{FC projection}}$ Draft context input $\xrightarrow{\text{diffusion}}$ 16 logits $\xrightarrow{\text{verify}}$ accepted tokens
\end{frame}

% ============================================================
% Slide 6 -- Code: Two-Phase KV Cache Write
% ============================================================
\section{Technical Challenges}

\begin{frame}[fragile]{Code: Two-Phase KV Cache Write}
\begin{lstlisting}
# Phase A: write padded context (with zeros) at cache_len
k_ctx_4d = k_ctx.transpose(1, 0, 2)[jnp.newaxis, :, :, :]
v_ctx_4d = v_ctx.transpose(1, 0, 2)[jnp.newaxis, :, :, :]
kv_cache_k = lax.dynamic_update_slice(
    kv_cache_k, k_ctx_4d, (0, 0, cache_len, 0)
)
kv_cache_v = lax.dynamic_update_slice(
    kv_cache_v, v_ctx_4d, (0, 0, cache_len, 0)
)

# Phase B: write noise at cache_len + actual_ctx_count,
# overwriting padding zeros from Phase A
noise_start = cache_len + actual_ctx_count
k_noise_4d = k_noise.transpose(1, 0, 2)[jnp.newaxis, :, :, :]
v_noise_4d = v_noise.transpose(1, 0, 2)[jnp.newaxis, :, :, :]
kv_cache_k = lax.dynamic_update_slice(
    kv_cache_k, k_noise_4d, (0, 0, noise_start, 0)
)
kv_cache_v = lax.dynamic_update_slice(
    kv_cache_v, v_noise_4d, (0, 0, noise_start, 0)
)
\end{lstlisting}
\vspace{-0.5em}
{\footnotesize Source: \texttt{models/jax/dflash.py} --- context+noise must be written in two phases because context length varies per step.}
\end{frame}

% ============================================================
% Slide 7 -- The seq_len Inflation Bug
% ============================================================
\begin{frame}[fragile]{The \texttt{seq\_len} Inflation Bug}
\textbf{Symptom:} $\tau = 2.49$, 10\% acceptance, 13 A/B tests inconclusive

\vspace{0.5em}
\textbf{Root cause:} \texttt{attn\_metadata.seq\_lens} includes \emph{unverified} draft tokens
\begin{itemize}
  \item Actual generated: 100 tokens
  \item Value passed to proposer: $100 + 16 = 116$ (includes draft tokens being verified)
  \item DFlash builds KV cache at wrong position $\Rightarrow$ stale context accumulates
\end{itemize}

\begin{alertblock}{The 4-line fix}
\begin{lstlisting}[backgroundcolor=\color{white},frame=none,xleftmargin=0pt]
accepted_seq_lens = self.runner.input_batch.num_tokens_no_spec[
    :attn_metadata.seq_lens.shape[0]].copy()
accepted_attn_metadata = replace(
    attn_metadata, seq_lens=device_array(
        self.runner.mesh, accepted_seq_lens.astype(np.int32)))
\end{lstlisting}
\end{alertblock}

\textbf{Result:} $\tau$: $2.49 \rightarrow 4.48$ \quad Speedup: $1.30\times \rightarrow 2.31\times$
\end{frame}

% ============================================================
% Slide 8 -- Why the Bug Was Hard to Find
% ============================================================
\begin{frame}{Why the Bug Was Hard to Find}
\begin{enumerate}
  \item \textbf{No visible errors} --- tokens were valid English, just suboptimal; model ``recovered'' from bad context
  \item \textbf{Eagle3 unaffected} --- Eagle3 has no persistent draft-side KV cache, so the same seq\_len value worked fine
  \item \textbf{Gradual accumulation} --- 10--16 phantom tokens added per decoding step; quality degrades slowly over generations
  \item \textbf{Bug is upstream} --- located in the speculative decoding \emph{manager}, not the DFlash proposer; proposer code was correct
\end{enumerate}
\end{frame}

% ============================================================
% Slide 9 -- Results: Standalone TPU vs GPU
% ============================================================
\section{Results}

\begin{frame}{Results: Standalone TPU vs GPU}
\centering
\footnotesize
\setlength{\tabcolsep}{3.5pt}
\begin{tabular}{l ccc ccc}
\toprule
& \multicolumn{3}{c}{\textbf{TPU v4}} & \multicolumn{3}{c}{\textbf{GPU (A100)}} \\
\cmidrule(lr){2-4} \cmidrule(lr){5-7}
\textbf{Dataset} & $\tau$ & Speedup & TPS & $\tau$ & Speedup & TPS \\
\midrule
GSM8K   & 6.25 & 3.67$\times$ & 222 & 6.69 & 5.53$\times$ & --- \\
\rowcolor{rowgreen}
Math500 & \textbf{8.72} & 4.49$\times$ & 276 & 7.84 & 5.12$\times$ & --- \\
AIME24  & 5.17 & 3.00$\times$ & 191 & 6.45 & 5.07$\times$ & --- \\
AIME25  & 6.53 & 3.53$\times$ & 218 & 7.30 & 5.38$\times$ & --- \\
\midrule
\textbf{Avg} & \textbf{6.67} & \textbf{3.67$\times$} & \textbf{227} & \textbf{7.07} & \textbf{5.53$\times$} & --- \\
\bottomrule
\end{tabular}

\vspace{0.8em}
\begin{itemize}
  \item TPU achieves \textbf{94\% of GPU $\tau$} (6.67 vs 7.07) --- algorithmic quality transfers
  \item Math500: TPU $\tau$ \textbf{exceeds} GPU (8.72 vs 7.84)
  \item Speedup gap ($3.67\times$ vs $5.53\times$) is hardware overhead, not algorithm quality
\end{itemize}
\end{frame}

% ============================================================
% Slide 10 -- Results: vLLM Pipeline
% ============================================================
\begin{frame}{Results: vLLM Pipeline}
\begin{columns}[T]
\begin{column}{0.45\textwidth}
\textbf{Pipeline Metrics (GSM8K)}
\vspace{0.3em}

\footnotesize
\begin{tabular}{lr}
\toprule
$\tau$ & 4.48 \\
Speedup & 2.31$\times$ \\
Spec TPS & 212 \\
Baseline TPS & 92 \\
Acceptance & 23.2\% \\
\bottomrule
\end{tabular}
\end{column}
\begin{column}{0.52\textwidth}
\textbf{Standalone vs Pipeline}
\vspace{0.3em}

\footnotesize
\begin{tabular}{lcc}
\toprule
& Standalone & vLLM \\
\midrule
$\tau$ & 6.25 & 4.48 \\
Speedup & 3.67$\times$ & 2.31$\times$ \\
TPS & 222 & 212 \\
\bottomrule
\end{tabular}

\vspace{0.5em}
\footnotesize
$\tau$ gap from vLLM scheduler/batch overhead (draft tokens rejected at batch boundaries). Same pattern reported in GPU DFlash paper.
\end{column}
\end{columns}
\end{frame}

% ============================================================
% Slide 11 -- DFlash vs Eagle3 on TPU
% ============================================================
\begin{frame}{DFlash vs Eagle3 on TPU}
\centering
\begin{tabular}{lcc}
\toprule
\textbf{Metric} & \textbf{DFlash} & \textbf{Eagle3} \\
\midrule
Draft tokens/step & 15 & 3 \\
$\tau$ & 4.48 & 2.18 \\
TPS & 212 & 78 \\
Speedup & 2.31$\times$ & 1.53$\times$ \\
\bottomrule
\end{tabular}

\vspace{1em}
\begin{itemize}
  \item DFlash $\tau$ is \textbf{2.06$\times$ Eagle3} on TPU (4.48 vs 2.18)
  \item Advantage transfers proportionally from GPU results
  \item Both measured under identical vLLM pipeline conditions
\end{itemize}
\end{frame}

% ============================================================
% Slide 12 -- Output Quality Verification
% ============================================================
\begin{frame}{Output Quality Verification}
\centering
\begin{tabular}{lc}
\toprule
\textbf{Comparison} & \textbf{Count (of 8)} \\
\midrule
Bit-exact match & 4 \\
Same final answer & 7 \\
Different answer & 1 \\
\bottomrule
\end{tabular}

\vspace{1em}
\begin{itemize}
  \item Divergences caused by \textbf{bf16 precision}: top-2 logits differ by $< 10^{-4}$; sampling flips at decision boundaries
  \item \textbf{Not} algorithmic error --- verified by inspecting logit distributions
  \item Divergent outputs are coherent alternative phrasings, not corruption
  \item Quality consistent with GPU speculative decoding guarantees
\end{itemize}
\end{frame}

% ============================================================
% Slide 13 -- Impact: The 4-Line Fix
% ============================================================
\begin{frame}{Impact: The 4-Line Fix}
\centering
\begin{tabular}{lccc}
\toprule
\textbf{Metric} & \textbf{Before} & \textbf{After} & \textbf{Change} \\
\midrule
$\tau$ & 2.49 & 4.48 & \textbf{+80\%} \\
Speedup & 1.30$\times$ & 2.31$\times$ & \textbf{+77\%} \\
Acceptance rate & 9.9\% & 23.2\% & \textbf{+134\%} \\
Tokens/sec & 120.8 & 212.4 & \textbf{+76\%} \\
\bottomrule
\end{tabular}

\vspace{1em}
\footnotesize
\begin{tabular}{l cccc}
\toprule
& \multicolumn{2}{c}{\textbf{$\tau$}} & \multicolumn{2}{c}{\textbf{Speedup}} \\
\cmidrule(lr){2-3} \cmidrule(lr){4-5}
\textbf{Dataset} & Before & After & Before & After \\
\midrule
GSM8K   & 2.49 & 4.48 & 1.30$\times$ & 2.31$\times$ \\
Math500 & 2.61 & 4.71 & 1.36$\times$ & 2.42$\times$ \\
AIME24  & 2.33 & 4.20 & 1.22$\times$ & 2.15$\times$ \\
AIME25  & 2.42 & 4.39 & 1.27$\times$ & 2.27$\times$ \\
\bottomrule
\end{tabular}
\end{frame}

% ============================================================
% Slide 14 -- Future Work
% ============================================================
\section{Future Work}

\begin{frame}{Future Work}
\begin{itemize}
  \item \textbf{PR to tpu-inference} --- upstream the DFlash proposer and seq\_len fix
  \item \textbf{TPU v5 benchmarks} --- v4 is experimental; v5 is fully supported with better matmul throughput
  \item \textbf{Concurrent serving benchmarks} --- 1--32 simultaneous clients; measure throughput under contention
  \item \textbf{Diffusion targets on TPU} --- port FailFast (diffusion-based target model) for end-to-end diffusion pipeline
\end{itemize}
\end{frame}

% ============================================================
% Slide 15 -- Summary
% ============================================================
\section{Conclusion}

\begin{frame}{Summary}
\begin{block}{What We Built}
Complete port of DFlash speculative decoding from GPU/PyTorch to TPU/JAX with full vLLM integration and zero upstream changes.
\end{block}

\begin{block}{Key Numbers}
\begin{itemize}
  \item Standalone: $\tau = 6.67$ (94\% of GPU), Math500 \emph{exceeds} GPU
  \item vLLM pipeline: $2.31\times$ speedup, 212 TPS
  \item DFlash $>$ Eagle3: $2.06\times$ higher $\tau$ on TPU
  \item Output quality: 7/8 same answer, 4/8 bit-exact
\end{itemize}
\end{block}

\begin{alertblock}{Key Finding}
A 4-line \texttt{seq\_len} fix doubled all performance metrics --- the hardest bugs leave no trace.
\end{alertblock}
\end{frame}

\end{document}
